\section{Alternatives and their evidence}\label{sec:alternatives}
\subsection{Mindfulness and acceptance-based treatments}
Mindfulness and acceptance-based approaches train attention and responses to thoughts and feelings rather than relying exclusively on disputing thought content.
Acceptance-based work overlaps with the broader CBT family \citep{hayes2006act}.
Its inclusion here does not establish an absence of cognitive or executive demands.
Brewer frames anxiety as a habit loop (trigger, worry, a feeling of control) and proposes noticing, curiosity about the reward value of worry, and a ``bigger, better offer'' \citep{brewer2021unwinding}.
Evidence for the specific programme is early: in a randomised trial of an app-based programme for GAD ($n=65$), anxiety scores fell 67\% versus 14\% under usual care, with mediation implicating reduced worry \citep{roy2021app}.
The trial is small, self-reported, compared with usual care rather than CBT, and has the programme's developer among its authors.

The following syntheses do not establish that mindfulness and acceptance approaches outperform evidence-based comparators.
A meta-analysis found essentially no difference ($g\approx0.00$) between mindfulness-based interventions and evidence-based treatments, often CBT \citep{goldberg2018mbi}; another found no difference from traditional CBT or behavioural therapies ($g=-.07$, nine comparisons; \citealt{khoury2013mindfulness}).
For acceptance and commitment therapy (ACT) versus CBT, the pooled difference was non-significant in one meta-analysis \citep{atjak2015act}, and a sceptical reading of 60 RCTs concluded that ACT is not yet well-established for any disorder \citep{ost2014act}.
In 276 adults with anxiety disorders, eight weeks of mindfulness-based stress reduction was noninferior to escitalopram, with fewer discontinuations for adverse events \citep{hoge2023mbsr}.
That is a non-inferiority result against a first-line drug, not evidence of superiority to CBT.
These comparisons do not establish universal equivalence across disorders, outcomes and follow-up periods.
Similar outcomes can arise through shared or different processes; they do not identify a mechanism by themselves.

\paragraph{Reward revaluation as a hypothesis}
An alternative self-regulation account proposes that attending to the experienced consequences of an action can change its reward value and reduce reliance on effortful restraint \citep{ludwig2020selfregulation}.
This is a mechanistic proposal, not comparative evidence against CBT.
An app-based eating study linked model-estimated reward values to reported eating, but had no control group \citep{taylor2021awareness}.
Its larger naturalistic sample showed decreasing estimated reward values without a corresponding decrease in reported eating behaviour across tool uses.
The study supports further testing of reward revaluation, not a causal conclusion about the intervention, a general mechanism for anxiety, or superiority to cognitive restructuring.

\subsection{Medication in ADHD}
In a network meta-analysis of 133 double-blind trials, amphetamines (standardised effect $-0.79$) and methylphenidate ($-0.49$) outperformed placebo in adults on clinician ratings, over short follow-up \citep{cortese2018adhdmeds}.
The certainty varied from high or moderate to low or very low.
A Cochrane review of extended-release methylphenidate found a self-rated effect of $-0.37$, judged 20 of 21 trials at high risk of bias, and noted that all trials excluded participants with psychiatric comorbidity such as anxiety or selected previous responders \citep{boesen2022ermph}.
So the pivotal medication trials say little about adults with ADHD \emph{and} anxiety.
